%% ============================================================
%%  PLANTILLA DE APUNTES DE MATEMÁTICA
%%  Compilar con:  pdflatex plantilla_apuntes.tex   (2 veces)
%%  Para el índice analítico: makeindex plantilla_apuntes
%%  y luego pdflatex otra vez. (Con --shell-escape se hace solo.)
%% ============================================================
\documentclass[11pt,a4paper,oneside]{report}

%% ---------- Idioma y codificación ----------
\usepackage[T1]{fontenc}
\usepackage[utf8]{inputenc}
\usepackage{lmodern}
\usepackage[spanish,es-noshorthands]{babel}
\usepackage{microtype}

%% ---------- Página y párrafos ----------
\usepackage[margin=2.5cm]{geometry}
\usepackage{parskip}            % separación entre párrafos sin sangría
\usepackage{setspace}
\setstretch{1.1}
\usepackage{fancyhdr}
\pagestyle{fancy}
\fancyhf{}
\fancyhead[L]{\small\nouppercase{\leftmark}}
\fancyhead[R]{\small\thepage}
\renewcommand{\headrulewidth}{0.4pt}

%% ---------- Matemática ----------
\usepackage{amsmath,amssymb,amsthm,mathtools}
\usepackage{bm}

%% ---------- Figuras, tablas, colores ----------
\usepackage{graphicx}
\graphicspath{{figuras/}}       % pon tus imágenes en la carpeta "figuras"
\usepackage{xcolor}
\usepackage{booktabs}           % tablas bonitas (\toprule, \midrule)
\usepackage{array}
\usepackage{tabularx}
\usepackage{longtable}
\usepackage{caption}
\usepackage{float}              % permite [H] para fijar posición
\usepackage{mwe}                % solo para imágenes de ejemplo; bórralo luego
\usepackage{tikz}

%% ---------- Código ----------
\usepackage{listings}
\lstset{
  basicstyle=\ttfamily\small,
  keywordstyle=\color{blue!70!black}\bfseries,
  commentstyle=\color{green!40!black}\itshape,
  stringstyle=\color{orange!80!black},
  numbers=left, numberstyle=\tiny\color{gray}, numbersep=8pt,
  backgroundcolor=\color{gray!8},
  frame=single, rulecolor=\color{gray!40},
  breaklines=true, showstringspaces=false, tabsize=4,
  literate={á}{{\'a}}1 {é}{{\'e}}1 {í}{{\'i}}1 {ó}{{\'o}}1 {ú}{{\'u}}1
           {Á}{{\'A}}1 {É}{{\'E}}1 {Í}{{\'I}}1 {Ó}{{\'O}}1 {Ú}{{\'U}}1
           {ñ}{{\~n}}1 {Ñ}{{\~N}}1
}

%% ---------- Índices ----------
\usepackage{imakeidx}
\makeindex[title=Índice analítico, intoc]

%% ---------- Cajas de colores ----------
\usepackage[most]{tcolorbox}

%% ---------- Hipervínculos y referencias (SIEMPRE al final) ----------
\usepackage[colorlinks=true,linkcolor=blue!60!black,citecolor=green!40!black,
            urlcolor=blue!60!black]{hyperref}
\usepackage[spanish,nameinlink]{cleveref}

%% ============================================================
%%  ENTORNOS: definición, teorema, demostración, notas...
%%  Todos comparten UN solo contador, dentro de cada capítulo.
%% ============================================================
\theoremstyle{definition}
\newtheorem{definicion}{Definición}[chapter]
\newtheorem{ejemplo}[definicion]{Ejemplo}
\newtheorem{ejercicio}[definicion]{Ejercicio}

\theoremstyle{plain}
\newtheorem{teorema}[definicion]{Teorema}
\newtheorem{proposicion}[definicion]{Proposición}
\newtheorem{lema}[definicion]{Lema}
\newtheorem{corolario}[definicion]{Corolario}

\theoremstyle{remark}
\newtheorem*{observacion}{Observación}
\newtheorem*{notademo}{Nota sobre la demostración}
\newtheorem*{idea}{Idea clave}

\renewcommand{\proofname}{Demostración}
\renewcommand{\qedsymbol}{$\blacksquare$}

% Nombres para cleveref
\crefname{definicion}{definición}{definiciones}
\Crefname{definicion}{Definición}{Definiciones}
\crefname{teorema}{teorema}{teoremas}
\Crefname{teorema}{Teorema}{Teoremas}
\crefname{proposicion}{proposición}{proposiciones}
\Crefname{proposicion}{Proposición}{Proposiciones}
\crefname{lema}{lema}{lemas}
\Crefname{lema}{Lema}{Lemas}
\crefname{corolario}{corolario}{corolarios}
\Crefname{corolario}{Corolario}{Corolarios}
\crefname{ejemplo}{ejemplo}{ejemplos}
\Crefname{ejemplo}{Ejemplo}{Ejemplos}
\crefname{ejercicio}{ejercicio}{ejercicios}
\Crefname{ejercicio}{Ejercicio}{Ejercicios}
\crefname{equation}{ecuación}{ecuaciones}
\Crefname{equation}{Ecuación}{Ecuaciones}
\crefname{figure}{figura}{figuras}
\Crefname{figure}{Figura}{Figuras}
\crefname{table}{tabla}{tablas}
\Crefname{table}{Tabla}{Tablas}
\crefname{chapter}{capítulo}{capítulos}
\Crefname{chapter}{Capítulo}{Capítulos}
\crefname{section}{sección}{secciones}
\Crefname{section}{Sección}{Secciones}

% Estilo visual: barra lateral de color + fondo suave
\tcbset{
  cajabase/.style={enhanced, breakable, sharp corners, boxrule=0pt,
    frame hidden, left=8pt, right=6pt, top=4pt, bottom=4pt,
    before skip=10pt, after skip=10pt}
}
\tcolorboxenvironment{definicion}{cajabase,
  colback=blue!6,   borderline west={3pt}{0pt}{blue!65!black}}
\tcolorboxenvironment{teorema}{cajabase,
  colback=red!6,    borderline west={3pt}{0pt}{red!70!black}}
\tcolorboxenvironment{proposicion}{cajabase,
  colback=orange!7, borderline west={3pt}{0pt}{orange!80!black}}
\tcolorboxenvironment{lema}{cajabase,
  colback=orange!7, borderline west={3pt}{0pt}{orange!80!black}}
\tcolorboxenvironment{corolario}{cajabase,
  colback=orange!7, borderline west={3pt}{0pt}{orange!80!black}}
\tcolorboxenvironment{ejemplo}{cajabase,
  colback=green!5,  borderline west={3pt}{0pt}{green!45!black}}
\tcolorboxenvironment{ejercicio}{cajabase,
  colback=green!5,  borderline west={3pt}{0pt}{green!45!black}}
\tcolorboxenvironment{notademo}{cajabase,
  colback=yellow!14, borderline west={3pt}{0pt}{yellow!60!black}}
\tcolorboxenvironment{observacion}{cajabase,
  colback=gray!10,  borderline west={3pt}{0pt}{gray!70}}
\tcolorboxenvironment{idea}{cajabase,
  colback=purple!6, borderline west={3pt}{0pt}{purple!70}}

%% ============================================================
%%  ATAJOS MATEMÁTICOS
%% ============================================================
\newcommand{\N}{\mathbb{N}}
\newcommand{\Z}{\mathbb{Z}}
\newcommand{\Q}{\mathbb{Q}}
\newcommand{\R}{\mathbb{R}}
\newcommand{\C}{\mathbb{C}}
\newcommand{\eps}{\varepsilon}
\newcommand{\dd}{\mathrm{d}}
\newcommand{\dpar}[2]{\frac{\partial #1}{\partial #2}}
\DeclarePairedDelimiter{\abs}{\lvert}{\rvert}
\DeclarePairedDelimiter{\norm}{\lVert}{\rVert}
\DeclarePairedDelimiter{\pesc}{\langle}{\rangle}
\DeclarePairedDelimiter{\conj}{\{}{\}}
\DeclareMathOperator{\sen}{sen}
\DeclareMathOperator{\tg}{tg}
\DeclareMathOperator{\Img}{Im}
\DeclareMathOperator{\Ker}{Ker}
\DeclareMathOperator{\rg}{rg}

% Atajo para "término indexado": lo escribe en el texto Y lo añade al índice
\newcommand{\termino}[1]{\textbf{#1}\index{#1}}

%% ============================================================
%%  PORTADA
%% ============================================================
\title{Apuntes de Topología}
\author{Práxedis Jiménez Ruvalcaba}
\date{\today}

\begin{document}

\maketitle
\tableofcontents

%% ============================================================
\part{Topología 1}
%% ============================================================

\chapter{Límites y continuidad}

\section{Definiciones y teoremas}

Los párrafos se separan con una línea en blanco. Para forzar un salto de línea
sin cambiar de párrafo, usa \verb|\\|:\\
esta línea queda justo debajo de la anterior.

\begin{definicion}[Límite de una sucesión]\label{def:limite}
  Sea $(a_n)_{n\in\N}\subset\R$. Decimos que $a_n$ \termino{converge} a
  $L\in\R$ si
  \[
    \forall \eps>0\;\exists N\in\N\;\text{tal que}\;
    n\ge N \implies \abs{a_n-L}<\eps .
  \]
\end{definicion}

\begin{teorema}[Unicidad del límite]\label{thm:unicidad}
  Si una sucesión converge, su límite es único.
\end{teorema}

\begin{notademo}
  La idea es suponer dos límites distintos y usar la desigualdad triangular
  con $\eps = \abs{L-L'}/2$ para llegar a una contradicción.
  Esta técnica se reutiliza en el \cref{cor:sumas}.
\end{notademo}

\begin{proof}
  Sean $L\neq L'$ dos límites de $(a_n)$ y tomemos $\eps=\tfrac{1}{2}\abs{L-L'}>0$.
  Por la \cref{def:limite}, existen $N_1,N_2$ tales que
  $\abs{a_n-L}<\eps$ si $n\ge N_1$ y $\abs{a_n-L'}<\eps$ si $n\ge N_2$.
  Para $n\ge\max\{N_1,N_2\}$:
  \begin{align}
    \abs{L-L'} &\le \abs{L-a_n}+\abs{a_n-L'} \label{eq:triangular}\\
               &< 2\eps = \abs{L-L'},
  \end{align}
  lo cual es absurdo.
\end{proof}

\begin{corolario}\label{cor:sumas}
  Si $a_n\to L$ y $b_n\to M$, entonces $a_n+b_n\to L+M$.
\end{corolario}

\begin{observacion}
  La desigualdad \eqref{eq:triangular} también aparece en la prueba del
  \cref{thm:unicidad}. Referencias automáticas: \cref{def:limite},
  \cref{thm:unicidad}, \cref{cor:sumas}; en el PDF cada una es un enlace
  clicable. Con \verb|\Cref| se escribe con mayúscula inicial al comienzo
  de frase, y con \verb|\cpageref| podrías citar la página.
\end{observacion}

\begin{ejemplo}
  La sucesión $a_n=1/n$ converge a $0$: dado $\eps>0$ basta tomar
  $N>1/\eps$.
\end{ejemplo}

\section{Código}

\begin{lstlisting}[language=Python, caption={Primeros términos de $1/n$}]
# Calcula los primeros términos de la sucesión
for n in range(1, 6):
    print(n, 1 / n)
\end{lstlisting}

En línea: la función \lstinline[language=Python]|math.sqrt(2)| devuelve $\sqrt2$.
Para código desde un archivo externo:
\verb|\lstinputlisting[language=Python]{codigo/ejemplo.py}|.

\section{Tablas}

\subsection{Tabla horizontal (encabezado en la primera fila)}

\begin{table}[H]
  \centering
  \caption{Valores de $f(x)=x^2$.}
  \label{tab:horizontal}
  \begin{tabular}{@{}lccccc@{}}
    \toprule
    $x$    & $-2$ & $-1$ & $0$ & $1$ & $2$ \\
    \midrule
    $f(x)$ & $4$  & $1$  & $0$ & $1$ & $4$ \\
    \bottomrule
  \end{tabular}
\end{table}

\subsection{Tabla vertical (encabezado en la primera columna)}

\begin{table}[H]
  \centering
  \caption{Propiedades de los conjuntos numéricos.}
  \label{tab:vertical}
  \begin{tabular}{@{}l >{\centering\arraybackslash}p{2.2cm}
                    >{\centering\arraybackslash}p{2.2cm}
                    >{\centering\arraybackslash}p{2.2cm}@{}}
    \toprule
    \textbf{Propiedad} & $\Q$ & $\R$ & $\C$ \\
    \midrule
    Cuerpo           & Sí & Sí & Sí \\
    Ordenado         & Sí & Sí & No \\
    Completo         & No & Sí & Sí \\
    \bottomrule
  \end{tabular}
\end{table}

Tabla con texto largo que se ajusta al ancho (\verb|tabularx|):

\begin{table}[H]
  \centering
  \begin{tabularx}{\textwidth}{@{}lX@{}}
    \toprule
    \textbf{Concepto} & \textbf{Descripción} \\
    \midrule
    Cota superior & Número mayor o igual que todos los elementos del conjunto,
                    que puede o no pertenecer a él. \\
    Supremo       & La menor de las cotas superiores. \\
    \bottomrule
  \end{tabularx}
\end{table}

Las referencias funcionan igual: ver \cref{tab:horizontal} y \cref{tab:vertical}.
Para tablas muy largas que cruzan páginas usa \verb|longtable|.

\section{Imágenes}

\begin{figure}[H]
  \centering
  % Reemplaza example-image-a por tu archivo, p. ej. grafico.png
  % (colócalo en la carpeta "figuras")
  \includegraphics[width=0.45\textwidth]{example-image-a}
  \caption{Pie de figura.}
  \label{fig:ejemplo}
\end{figure}

\Cref{fig:ejemplo} muestra una imagen. También puedes dibujar con TikZ:

\begin{figure}[H]
  \centering
  \begin{tikzpicture}[scale=1.2]
    \draw[->] (-0.3,0) -- (3,0) node[right] {$x$};
    \draw[->] (0,-0.3) -- (0,2.5) node[above] {$y$};
    \draw[domain=0:1.55, smooth, thick, blue] plot (\x, {\x*\x});
  \end{tikzpicture}
  \caption{Gráfico de $y=x^2$ dibujado con TikZ.}
  \label{fig:tikz}
\end{figure}

%% ------------------------------------------------------------
\chapter{Segundo tema del curso}

\begin{teorema}[Ejemplo de cita cruzada]\label{thm:cruzado}
  Usando el \cref{thm:unicidad} del \cref{chap:limites-ref}... (cita a otro
  capítulo con \verb|\cref{chap:...}| tras poner un \verb|\label| en él).
\end{teorema}

\chapter{Tema con etiqueta}\label{chap:limites-ref}

Texto del capítulo referenciado arriba. Recuerda: ejecuta \texttt{pdflatex}
dos veces para que las referencias se resuelvan.

%% ============================================================
\part{Nombre del segundo curso (p.\,ej. Álgebra Lineal)}
%% ============================================================

\chapter{Espacios vectoriales}

\begin{definicion}[Espacio vectorial]\label{def:ev}
  Un \termino{espacio vectorial} sobre un cuerpo $\mathbb{K}$ es un conjunto $V$
  con operaciones suma y producto por escalares que cumplen los axiomas usuales.
\end{definicion}

\begin{teorema}[Teorema de la dimensión]\label{thm:dim}
  Sea $T:V\to W$ lineal con $\dim V<\infty$. Entonces
  \[
    \dim V = \dim\Ker T + \dim\Img T .
  \]
\end{teorema}

\begin{proof}
  Sea $\{u_1,\dots,u_k\}$ una base de $\Ker T$, que se extiende a una base
  $\{u_1,\dots,u_k,v_1,\dots,v_m\}$ de $V$. Se prueba que
  $\{Tv_1,\dots,Tv_m\}$ es base de $\Img T$.
\end{proof}

\begin{idea}
  Extender una base del núcleo a una base de todo el espacio.
\end{idea}

Se usa el \cref{thm:dim} junto con la \cref{def:ev}. Recuerda además que la
matriz $A=\begin{pmatrix} 1 & 2\\ 3 & 4\end{pmatrix}$ tiene $\rg A=2$.

\part{Topología 2}
\chapter{Tercer Parcial}
\section{Examen}
\begin{ejercicio}
Sean $f_0\simeq f_1$ rel $\dot{I}$ caminos en $X$ y sean $g_0\simeq g_1$
rel $\dot{I}$ camino en $Y$. Si para cada $i \in {0,1}$, $(f_i,g_i)$ es el camino en $X$x$Y$ 
definidio por $t\rightarrow (f_i(t),g_i(t))$, demuestre que $(f_0,g_0) \simeq (f_1,g_1)$ rel $\dot{I}$.
\end{ejercicio}
como $f_0\simeq f_1$ rel $\dot{I}$ $\Rightarrow \exists F:I$x$I \rightarrow X$ donde $I=[0,1]$ y $F$ es 
una homotopía tal que $\forall s \in [0,1]$\\
$F(t,0)=f_0(s)$\\
$F(t,1)=f_1(s)$\\
y $F(0,s)=f_0(0)=f_1(0)$,$F(1,s)=f_0(1)=f_1(1)$ $\forall s \in [0,1]$\\
De manera análoga para $g_0 \simeq g_1$ rel $\dot{I}$ tendríamos una homotopía $G:I$x$I\rightarrow Y$. 
Con ambas definimos $H:I$x$I\rightarrow X$x$Y$ como $H(t,s)=(F(t,s),G(t,s))$.\\
$F$ y $G$ son continuas por ser homtopías, por teorema sobre producto de continuas visto en Topología 1, sabemos
entonces que $(F,G)$ es continua, es decir; $H$ es continua.\\
También podemos notar sobre los puntos extremos que\\
$H(t,0)=(F(t,0),G(t,0))=(f_0(t),g_0(t))$ $\forall s \in [0,1]$ y por tanto, 
$H(t,1)=(F(t,1),G(t,1))=(f_1(t),g_1(t))$. Lo cual ya probabría que $H$ es una homotopía. 
Pero también podemos notar que\\
$H(0,s)=(F(0,s),G(0,s))=(f_0(0),g_0(0))$ $\forall t \in [0,1]$, y que 
$H(1,s)=(F(1,s),G(1,s))=(f_1(1),g_1(1))$. Donde ya sabemos que $f_0(0)=f_1(0)$, $g_0(0)=g_1(0)$,
$f_1(1)=f_0(1)$, $g_1(1)=g_0(1)$. Entonces tendríamos la garantía de que $H$ es una homotopía en la cual\\
$(f_0,g_0) \simeq (f_1,g_1)$ rel $\dot{I}$.\\
Donde $(f_0(t),g_0(t))=(f_0,g_0)$ y $(f_0(t),g_0(t))=(f_0,g_0)$ rel $\dot{I} \blacksquare$

\begin{ejercicio}
  Sea $X$ el subespacio topológico de $\R^2$ cuyo conjunto base es la clausura del conjunto $S=\conj{\sen(\frac{1}{x})|x\in (0,1]}$
  sobre $\R^2$. Demuestre que $\pi (X,x_0)$ es isomorfo al grupo trivial para cada $x_0 \in X$
\end{ejercicio}

De topología 1 sabemos que $\bar{S}=S \cup (\conj{0},[-1,1])$. Denotemos entonces $L=(\conj{0},[-1,1])$ y veamos
que si $(0,y) \in L \Rightarrow \exists \theta_0$ tal que $\sen(\theta_0)=y$\\
Si $x_n=\frac{1}{\theta_0 +2n\pi}$ con $n\in\N$\\
$\Rightarrow x_n \in (0,1]$ y $\lim_{n\rightarrow \infty} x_n=0$\\
Si ahora definimos $P_n=(x_n,\sen(\frac{1}{x}))=(x_n,\theta_0 +2n\pi)=(x_n,y)\in S$ pero, \\
$\lim_{n\rightarrow \infty} P_n=(\lim_{n\rightarrow \infty} x_n, \lim_{n\rightarrow \infty} y)=(0,y)$\\
$\therefore L \subset S'$\\
Por otro lado,
si $(x_n,y_0)\in S'\Rightarrow \exists q_m=(x_m,\sen(\frac{1}{x_m}))$ tal que $q_m\in S$ y 
$q_m \rightarrow (x_0,y_0)$\\
como $\lim_{m\rightarrow \infty} x_m=x_0$  \\
Si $x_0>0 \Rightarrow \sen(\frac{1}{x_m})$ es continuo y por continuidad
$\lim_{m\rightarrow \infty} \sen(\frac{1}{x_m})=\sen(\lim_{m\rightarrow \infty}\frac{1}{x_m})=\sen(\frac{1}{x_0})$
entonces, $\lim_{m\rightarrow \infty}=(x_0,y_0)\in S$. \\
Si $x_0=0$ como $im(\sen(x))=[-1,1]$ $\Rightarrow \lim_{m\rightarrow \infty} \sen(\frac{1}{x_m}) \in[-1,1] \forall m$ \\
$\therefore (x_0,y_0)=(0,y)$ tal que $y\in[-1,1]$\\
$\therefore S' \subset L$\\
$\therefore L = S'$\\
Ahora revisemos que todo convexo es conexo por caminos\\
Si $X$ es convexo y $x_0,x_1 \in X$ $\Rightarrow (1-t)x_0 + tx_1 \in X$ $\forall t \in [0,1]$.
Si $\alpha(t)=(1-t)x_0 + tx_1$ $\Rightarrow \alpha(0)=x_0$ y $\alpha(1)=x_1$ $\Rightarrow \alpha(t) \subset X$
$\forall x_0,x_1 \in X$.\\
$\therefore$ Siempre existe camino continuo entre dos elementos arbitrarios de $X$ convexo. Es decir, si $X$ es convexo 
$\Rightarrow$ es conexo por caminos.\\
$\conj{0}$x$[-1,1]$ es un espacio convexo y por lo anterior, es conexo por caminos.\\
Revisemos que $S$ es también convexo, pues notemos que $(0,1) \simeq S$. Sea $f(x)=(x,\sen(\frac{1}{x}))$, $f:(0,1] \rightarrow S$.
Donde por $x$ y $\sen(\frac{1}{x})$ ser continuas en el dominio de $f$, entonces $f$ es continua por argumento
de producto de continuas. Por la propia definición de $S$, $f$ es sobre. Además, si $f(x)=f(x_0)$ 
$\Rightarrow (x,\sen(\frac{1}{x}))=(x_0,\sen(\frac{1}{x_0})) \iff x=x_0$, así que $f$ es también inyectiva.
Si además definimos $f^{-1}(x,\sen(\frac{1}{x}))=x$, la cual es una proyección de una función continua, que por teorema 
es tambiénn continua, así que $f^{-1}$ es también continua. \\
$\therefore f $ es homomorfismo $(0,1]\cong S$\\
Y como $(0,1]$ es convexo y esta es invariante topológica, entonces $S$ es también convexo $\Rightarrow$ es conexo por caminos.\\
Ahora notemos que $\bar{S}$ no es conexo por caminos, específicamente, no existe camino entre $S$ y $L$ continuo. 
Esto se puede argumentar tomando en cuenta que el límite de $\sen(x)$ no existe. Así que, no podemos contar con caminuos
continuos de $S$ a $L$. \\
Si suponemos que $\exists \gamma:[0,1]\rightarrow x$ tal que $\gamma(0)=(0,y_0)\in L$ y $y(1)=(x_1,y_1)\in S$,
donde $y$ es un camino continuo. Entonces, $y(t)=(x(t),y(t))$ nos garantiza que $x(t)$ y $y(t)$ son continuas.\\
Si $k=\conj{t|x(t)=0} \rightarrow x(k)={0}$ que es cerrado,
así que k es cerrado por continuidad de $x(t)$, lo cual a su vez implica que k tiene un máximo. Si $max K =t_i$,
entonces, $\forall t \in (t_i,1]$ $\gamma(t) \in S$. \\
Si construimos una sucesión $t_n\rightarrow t^+_i$ $\Rightarrow \gamma(t_n)\in S$. Pero $\lim_{n\rightarrow \infty^+}y(t_n)=y(t_i)$ 
que significa $\lim_{n\rightarrow \infty^+}\sen(\frac{1}{x(t_n)})$, pero ya habíamos mencionado que el límite de $\sen(x)$ no existe,
lo cual es una contradicción, así que $\gamma$ no puede ser continua. \\
$\therefore \bar{S}$ no es conexo por caminos. Pero $S$ y $L$ sí lo son, más aún, por teorema 2.34, por ser ambos 
conexos, entonces son contrácitles. \\
$\pi_1(X,x_0)$ son los caminos cerrados con base en $x_0$ donde la imagen de cualquier camino continuo debe estar 
contenida en una única componente por camino. Así que tenemos dos caos $x_0 \in S$ o $x_0 \in L$, ya que 
$\pi_1(\bar{S},x_0)\cong \pi_1(S,x_0)$ o $\pi_1(\bar{S},x_0)\cong \pi_1(L,x_0)$, dependiendo de si $x_0 \in S$ o 
dónde esté por el teorema 3.31. \\
Como $S$ es contráctil, por teorema 3.35.2 podemos argumentar que $\pi_1(S,x_0)\cong \conj{e}$ si $x_0 \in L$.
De manera análoga $\pi_1(L,x_0)\cong \conj{e}$.\\
$\therefore \forall x_0 \in \bar{S}$ $\pi_1(\bar{S},x_0)\cong \conj{e} \blacksquare$

\begin{ejercicio}
  Suponga que exitse un espacio topológico puntuado $(X,x_0)$, tal que $\pi_1 (X,x_0) \cong \Z$. Demuestre que para cada $n \in \Z$
  existe espacio topológico puntuado $(X_n, y_{n,0})$ tal que $\pi_1(Y_n, y_{n,0}) \cong \Z^n$
\end{ejercicio}

Por teorema 3.33 sabemos que dados $(X,x_0)$, $(Y,y_0)$ $\Rightarrow \pi_1(X \times Y,(x_0,y_0)) \cong \pi_1(X,x_0) \times \pi_1(Y,y_0)$\\
Sea $Y_n= X\times X \times ... \times X$ ($n$ productos de $X$), con $y_{(n,0)}=(x_0,...,x_0)$ ($n$-tupla de $x_0$). Entonces, \\
$\pi_1(Y_n,y_{(n,0)})=\pi_1(X\times X \times ... \times X,(x_0,...,x_0))$\\
$\cong \pi(X,x_0) \times\pi(X,x_0) \times... \times \pi(X,x_0)$\\
$\cong \Z \times \Z \times ... \times \Z$\\
$\cong \Z^n$ \\
Revisemos entonces por inducción que $\pi_1(Y_n, y_{n,0}) \cong \Z^n$ es una proposición válida. Notemos que el 
caso donde $n=1$ es la hipótesis del enunciado del problema. $n=2$ Está dado por el teorema 3.33 directamente, pues 
$\pi_1(X \times X, (x_0, x_0))\cong \pi_1 (X,x_0)\times  \pi_1 (X,x_0)\cong \Z \times \Z \cong \Z ^n$.\\
Supongamos que la proposición es válido para un valor admisible de la variable, digamos $k$\\
Hipótesis de inducción: $\pi_1(Y_k,y_{(k,0)})=\pi_1(X\times ... \times X,(x_0,...,x_0)) \cong \Z^k$.\\
Verifiquemos si la proposición es válida para el siguiente valor admisible, es decir, $k+1$\\
$Y_{k+1}=Y_k \times X$ $\Rightarrow \pi_1(Y_{k+1},y_{(k+1,0)})=\pi_1(Y_k\times X, (y_{(k,0)},x_0))$ que por teorema 3.33
$pi_1(Y_k\times X, (y_{(k,0)},x_0))\cong \pi_1(Y_k,y_{(k,0)}) \times \pi_1(X,x_0)$\\
$\cong \Z^k \times \Z$\\
$\cong \Z^{k+1}$\\
$\therefore \pi_1(Y_{k+1},y_{(k+1,0)})\cong \Z^{k+1}$\\
$\therefore$ La proposición es válida $\forall n \in \N$
$\therefore \pi_1(Y_n, y_{n,0}) \cong \Z^n$ $\forall n \in \N \blacksquare$

\begin{ejercicio}
  Sea $X$ un espacio topológica. Muestre que existe una categorí cuyos objetos son los elementos de $X$, sus morfismos son: \\
  $Hom(p,q)=\conj{[f] clase de caminos | \alpha [f]=p, \omega[f]=q}$\\ 
  y la composición \\
  $Hom(p,q)\times Hom(q,r) \rightarrow Hom(p,r)$\\
  definido por\\
  $([f],[g])\rightarrow [f*g]$\\
  Muestre que todo morfismo $C$ es una equivalencia.
\end{ejercicio}
La concatenación de caminos ya se probó en clase que es asociativa, por lo que en este caso solo hay que revisar 
la identidad para probar que $C$ es en efecto, una categoría.\\
Como $obC=X \Rightarrow \forall p \in X$ $I_p \in Hom(p,q)$ es la identidad de $p$ que sería $I_p=[e_p]$ con 
$ep:[0,1] \rightarrow X$ tal que $e_p(t)=p \forall t \in [0,1]$. Así que dado $[f]\in Hom(p,q)$ tenemos que 
$I_p*[f]=[f]$ y $[f]*I_q=[f]$ que por parcial 2 sabemos que esto significa que $e_p*f\simeq f$ y $f*e_q\simeq f$. 
Por ello se propone $N_i:IxI\rightarrow X$ donde:\\
$N_i(s,t) = \left\{
  \begin{array}{l}
    p \quad 0\leq s \leq t/2 \\
    f(\frac{2s-t}{2-t}) \quad t/2 \leq s \leq 1 \\
    f(\frac{2(t/2)-t}{2-t})=f(0)=p
  \end{array}
  t \in [0,1]
\right.$\\

$N_i(s,0) = \left\{
  \begin{array}{l}
    p=f(s) \quad 0\leq s \leq 0  \\
    f(s) \quad 0 \leq s \leq 1 
  \end{array}
\right.$\\

Que cabe remarcar, son continuas por pasting lemma, ya que $p$ es constnte y $f$ inicia en p, cuando $f$ es camino 
y por tanto, continua.

$N_i(s,1) = \left\{
  \begin{array}{l}
    p \quad 0\leq s \leq 1/2  \\
    f(\frac{2s-1}{2-1}) \quad 1/2 \leq s \leq 1 
  \end{array}
\right.$\\

Entonces $N_i(s,1)= I_p*f(s)$, así que $N_i$ va de $f$ a $I_p*f(s)$ y $N_i(0,t)=p$ por $s=0$. Mientras que $N_i(1,t)=f(\frac{2-t}{2-t})=q$\\
$\therefore$ Los extremos $p$ y $q$ se mantienen\\
$\therefore N_i$ es homotopía tal que $f\simeq e_p*f$. \\
De manera análoga establecemos $N_d:[IxI]\rightarrow X$
como:\\

$N_d(s,t) = \left\{
  \begin{array}{l}
    f(\frac{2s}{2-t}) \quad 0\leq s \leq 1- t/2 \\
    q \quad 1-t/2 \leq s \leq 1 \\
    f(\frac{2(1-t/2)-t}{2-t})=f(1)=q
  \end{array}
\right.$\\
$N_d(0,t)=f(0)=p$\\
$N_d(1,t)=q$\\

$N_d(s,0) = \left\{
  \begin{array}{l}
    f(s) \quad 0\leq s \leq 1  \\
    q \quad 1 
  \end{array}
\right.$\\

$N_i(s,1) = \left\{
  \begin{array}{l}
    f(s) \quad 0\leq s \leq 0  \\
    q \quad 1/2 \leq s \leq 1 
  \end{array}
\right.$\\
Y por los argumentos del punto anterior tendríamos que $N_d$ es una homotopía tal que $f(s)*e_q\simeq f(s)$\\
$\therefore I_p$ e $I_q$ son los neutros de los morfismos.\\
$\therefore C$ es una categoría.\\
Ahora para revisar que todo morfismo en $C$ es una equivalencia, debemos hallar los inversos $\forall [f] \in Hom$.
Así que tomemos $[f \in Hom(p,q)]$\\
Deberíamos proceder de manera similar, hallando las homotpías para poder probar que $[f]*[f^{-1}]=I_p$ y $[f^{-1}]*[f]=I_q$\\
Sea $f\in[f]\Rightarrow f:I\rightarrow X$ es camino continuo tal que $f(0)=p$ y $f(1)=q$; $f^{-1}:I\rightarrow X$ lo
definimos como $f(1-t)$ la cual es continua al ser una reparametrización con mismo dominio y además $\alpha(f^{-1})=q$
y $w(f^{-1})=p$ entonces, $[f^{-1}]\in HOm(q,p)$. \\
Ahora definimos:\\
$H(s,t) = \left\{
  \begin{array}{l}
    f(2s(1-t)) \quad 0\leq s \leq 1/2  \\
    f(2(1-t)(1-s)) \quad 1/2 \leq s \leq 1 
  \end{array}
  t \in [0,1]
\right.$\\
donde nuevamente $H$ es continua por pasting lema y porque ambas reparametrizaciones se mantienen continuas. Notése 
que ambas son $f(1-t)$ en $s=1/2$. Veamos que\\

$H(s,0) = \left\{
  \begin{array}{l}
    f(2s) \quad 0\leq s \leq 1/2  \\
    f(2(1-s)) \quad 1/2 \leq s \leq 1 
  \end{array}
  (f*f^{-1})(s)
\right.$\\

$H(s,1) = \left\{
  \begin{array}{l}
    f(0) \quad 0\leq s \leq 1/2  \\
    f(0) \quad 1/2 \leq s \leq 1 
  \end{array}
  e_p(s)
\right.$\\

Así que $H$ va de $(f*f^{-1})(s)$ a $e_p(s)$ de manera continua y 
$H(0,t)=f(0)=p=f^{-1}(1)$ y $H(1,t)=f(0)$. Así que $\forall t\in [0,1]$ se mantienen los extremos fijos.\\
$\therefore H$ es una homotopía tal que $(f*f^{-1})(s)\simeq e_p(s)$\\
$\therefore[f]*[f^{-1}]=I_p$\\
De manera análoga podríamos revisar que \\
$G(s,t) = \left\{
  \begin{array}{l}
    f(1-2s(1-t)) \quad 0\leq s \leq 1/2  \\
    f(1-2(1-t)(1-s)) \quad 1/2 \leq s \leq 1 
  \end{array}
  t \in [0,1]
\right.$\\

Nos da la homotopía $f^{-1}*f\simeq e_q$. Observando que en $s=1/2$ $f(1-(1-t))=f(t)=f(1-2(1/2)(1-t))$, tendríamos 
que aplica el pasting lema para continuidad y que $G(0,t)=G(1,t)=f(1)=q$, $G(s,0)=(f^{-1}*f)(s)$ y $G(s,1)=e_q(s)$.\\
$\therefore f^{-1}*f \simeq e_q$\\
Así que $[f^{-1}]*[f]=I_q$ \\
$\therefore$ Como todo morfismo posee un inverso, $C$ es una equivalencia $\blacksquare$.
\begin{ejercicio}
  Demuestre que si existe algún espacio topológico puntuado $(Y,Y_0)$ que cumple con que $\pi_1(Y,Y_0)\not\cong \conj{e_Y}$
  $\Rightarrow \pi_1(S',1) \not\cong \conj{e_{S'}}$
\end{ejercicio}
Durante la demostración del teorema 3.30 vimos que si tomamos $h:(X,X_0) \rightarrow (Y,Y_0)$ continua, entonces 
$\pi_1(h)$ es un homomorfismo tal que $\pi_1(X,X_0)\rightarrow \pi_1(Y, Y_0)$ donde $h([\gamma])=[h\circ \gamma]$
con $[\gamma] \in \pi_1 (X,X_0)$. \\
Si tomamos $\alpha : [0,1]\rightarrow Y$ camino cerrado tal que $\alpha(0)=\alpha(1)=Y_0$. Y notamos que $S'$ es 
en realidad un cociente del [0,1] identificando el 0 con el 1, entonces $q(t)=(\cos(2\pi t), \sen(2 \pi t))$ es
un mapeo natural que va de $q:[0,1]\rightarrow S' \cong [0,1]/\conj{0,1}$, con esto retomamos el teorema 1.11 
donde $g:S'\rightarrow Y$ es continua $\iff \alpha = g \circ q$ es continua, y como es un camino cerrado tenemos
entonces que $g$ es continua.\\
Si $\alpha \cong Y_0$ por teorema 2.31, $g \cong Y_0$ y entonces tendríamos una biyección entre $\alpha$ y $g$. 
Si suponemos que $\pi_1(S',1) \cong \conj{e_{S'}}$, por hipótesis $\pi_1(Y,Y_0)\not\cong \conj{e_Y}$, así que
existe $[\alpha]in \pi_1(Y,Y_0)$ tal que $[\alpha]\not =I_Y$. Pero por continuidad de $g$, $\pi_1(g)$ es homomorfismo. 
Que si lo evalúamos en $q$, entonces $\pi_1(g)([q])=[g \circ q]=[\alpha]$, pero como supusimos que 
$\pi_1(S',1)\cong \conj{e_{S'}}$ y $\pi_1(g)$ es homomorfismo, entonces $\pi_1(g)([q])=\conj{e_Y}=I_Y$, entonces
$[\alpha]=I_Y$ lo cual es una contradicción.
$\therefore \pi_1(S', Y_0) \not \cong \conj{e_{S'}}$
%% ============================================================
%%  ÍNDICE ANALÍTICO
%% ============================================================
\printindex

\end{document}
